\chapter{ Suivi de projet }
\label{chap4}
\minitoc
%=====================================%
\section{Introduction}
%\label{sec11}
%=====================================% 
Le succès du projet Ingénieur Engagé repose autant sur la qualité technique de la solution que sur la méthode de gestion appliquée. Ce chapitre présente les outils, méthodes et pratiques adoptés pour piloter efficacement le projet, assurer la coordination entre les membres de l’équipe, garantir le respect des délais, et maintenir une relation constructive avec le client et les utilisateurs.
%=====================================%
\section{Répartition des rôles et implication de l’équipe}
%=====================================%
Le projet a été conduit par une équipe de trois membres :
\begin{itemize}
    \item Un chef de projet, chargé de l’organisation globale et du suivi technique.
    \item Une responsable communication, en charge des échanges avec le client, de la coordination fonctionnelle, de la mise en forme des livrables, ainsi que de la réalisation des audits.
    \item Un troisième membre
\end{itemize}
Cette structure simple a permis une bonne répartition des tâches et une progression fluide du projet.

\subsection{Détails des contributions au projet}
L’ensemble de l’équipe a contribué activement à la phase de conception, aussi bien sur les aspects fonctionnels que techniques. Cette étape a permis de poser des bases solides à travers la formalisation des besoins, la modélisation UML et la création des maquettes UI/UX.

Le projet comportait un nombre significatif de fonctionnalités, avec des cas d’usage variés à couvrir, plusieurs rôles utilisateurs à gérer, et des règles de gestion à implémenter.
Conscients de ces enjeux, nous avons fait tout notre possible pour maintenir une organisation rigoureuse, respecter les délais imposés, et adapter notre planification au fil de l’avancement.
Cette rigueur nous a permis de livrer chaque jalon dans les temps et d’aborder chaque étape du projet de façon structurée et efficace.

Deux propositions de maquettes d’interface ont été réalisées : l’une par Nabilou, l’autre par Ines. Elles ont été présentées au client, puis fusionnées dans la version finale validée.
Wiame a notamment participé à la conception des diagrammes UML (use cases, classes), ainsi qu’à la relecture et l’enrichissement du dossier de conception, en coordination avec l’équipe.

Lors du développement, la répartition des rôles s’est faite en réponse aux défis techniques et organisationnels.
Cette répartition structurée nous a permis de développer de nouvelles compétences, tout en assurant une couverture complète du projet sur les volets front-end, back-end et intégration.

Nabilou s’est concentré sur le développement back-end, en prenant en charge la logique applicative, l’API, la base de données, et le suivi technique.

Ines a réalisé l’intégralité du front-end, avec une attention particulière portée à l’ergonomie et la cohérence visuelle.

L’intégration entre le front-end et le back-end a été assurée de manière conjointe, impliquant la création ou l’ajustement de plusieurs fonctionnalités dans les deux couches, afin d'assurer une communication fluide et fiable entre l'interface utilisateur et le système.

Ines a également rédigé et mis en forme tous les livrables en LaTeX, assuré la mise en page du rapport final, et rédigé tous les comptes rendus de réunions tout au long du projet.
Enfin, Wiame a apporté un soutien actif dans les tests fonctionnels, la relecture, et la structuration des livrables, contribuant pleinement à la qualité finale du rendu.


\begin{table}[h]
\centering
\begin{tabular}{|>{\bfseries}m{4cm}|m{5cm}|>{\centering\arraybackslash}m{4cm}|}
    \hline
    Membre & \underline{Rôle} & Pourcentage de participation \\
    \hline
    Nabilou ANOIR & Chef de projet & 40\% \\
    \hline
    Inès GRIBAA & Responsable communication & 40\% \\
    \hline
    Wieme Bourjila & Membre & 20\% \\
    \hline
\end{tabular}
\vspace{0.3cm} % petit espace après le tableau pour aérer
\caption{Répartition des rôles et du pourcentage de participation}
\label{tab:repartition}
\end{table}


Nous avons opté pour une approche agile adaptée à notre contexte étudiant, en combinant des principes de Scrum et de Kanban :
\begin{itemize}
    \item Sprints courts de 1 à 2 semaines avec objectifs définis.
    \item ScrumBan board sur Trello : suivi des tâches selon leur état (à faire, en cours, terminé).
    \item Diagramme de Gantt : planning initial et ajustements réguliers selon l’avancement.
    \item Réunions hebdomadaires pour synchroniser les efforts et ajuster les priorités.
\end{itemize}


\begin{figure}[H]
	\begin{center}
		\fbox{ \includegraphics[scale=0.75]{logos/captTrello.jpg}}
		\caption{Tableau ScrumBan}
		\label{fig:11}
	\end{center}
\end{figure}

\subsection{ Communication et collaboration }
\begin{itemize}
     

    \item Utilisation de GitHub pour le versioning et les relectures croisées.
    \item Google Docs pour la rédaction collaborative (rapport, note de cadrage, comptes rendus).
    \item WhatsApp et Google Meet pour la communication rapide et les réunions distancielles.
\end{itemize}

Cette organisation a permis une excellente réactivité et une traçabilité complète du travail effectué.


\subsection{Implication du client }
\begin{itemize}
     

    \item Des entretiens ont été menés avec le client et les utilisateurs finaux pour identifier les besoins.
    \item Des formulaires ont été diffusés auprès des étudiants pour recueillir leurs attentes, accompagnés d’entretiens menés avec les référents et plusieurs étudiants pour affiner la compréhension des besoins.
    \item Les maquettes Figma ont été validées étape par étape avec le client.
\end{itemize}


\subsection{ Retours utilisateurs}
Tout au long du projet, l'implication des utilisateurs finaux a constitué un axe essentiel pour garantir l’adéquation de la solution avec les besoins réels du terrain. Dès les premières phases de cadrage, des échanges réguliers ont été menés avec les étudiants, les référents et la direction afin de recueillir leurs attentes, valider les choix fonctionnels, et tester les premières maquettes.
Ces retours ont permis d’ajuster certaines fonctionnalités clés de l’application, notamment :
\begin{itemize}
    \item  La simplification du processus de déclaration d’actions pour les étudiants,

    \item L’ajout de filtres et de statuts explicites dans le tableau de bord,
    
    \item L’optimisation du processus de validation par les référents,
    
    \item La centralisation des actions validées dans une interface claire pour la direction et la scolarité.
\end{itemize}
Afin de structurer et formaliser la collecte de retours sur la version finale de l’application (MVP), un audit utilisateur a été mis en place. Il s’agit d’un formulaire Google Forms diffusé auprès des différentes catégories d’utilisateurs (étudiants, référents, direction, scolarité), visant à évaluer la facilité d'utilisation, la clarté de l’interface, la pertinence des fonctionnalités proposées, et recueillir des suggestions d’amélioration.
Le formulaire comporte à la fois :
\begin{itemize}
    \item des questions fermées (échelle de satisfaction, oui/non, menus déroulants),

     \item des questions ouvertes permettant aux participants de s’exprimer librement.
\end{itemize}

\href{https://docs.google.com/forms/d/e/1FAIpQLSfiW_80BxFX8UwUl4m0hySDFZ3xID2WkFncai5-XlAUqB_NmQ/viewform?usp=header}{Lien vers le formulaire d’audit utilisateur}



À la date de rédaction de ce rapport, la collecte des réponses était encore en cours. L’analyse des résultats sera réalisée dans une phase ultérieure, en vue d’alimenter une roadmap d’amélioration continue et de préparer une version enrichie de la plateforme. Ce retour d’expérience sera déterminant pour prioriser les prochaines évolutions de l’application selon les usages concrets et les retours du terrain.






  \subsubsection{ Gestion des priorités, qualité et livrables}
\begin{itemize}
     \item Les fonctionnalités ont été priorisées selon leur criticité (authentification, validation des actions, attribution des points).
    \item Mise en place de tests fonctionnels et vérification de chaque étape via GitHub.
    \item Recettage réalisé avant livraison pour garantir la conformité du produit final.
    \item Documentation complète fournie avec l’application.
\end{itemize}
 







\section{Conclusion}
 Le pilotage du projet a été rigoureux, structuré et collaboratif. Grâce à une répartition claire des rôles, des outils adaptés et une communication fluide, nous avons pu respecter les délais, assurer la qualité des livrables, et livrer une solution alignée sur les attentes du client. Ce suivi de projet démontre notre capacité à travailler en équipe de manière professionnelle et agile.

